% Blog post template.
% Copy this file to blog/posts/<slug>.tex. The file name becomes the URL
% (/blog/<slug>.html). Commit it and the pre-commit hook builds the page and feed.
% It still compiles with pdflatex, so you can preview it as a PDF too.
%
% Works: sections, lists, emphasis, amsmath environments, \newcommand macros,
%   \label/\eqref for equations, footnotes, tables, \cite with \bibliography.
% Doesn't: numbered \ref to sections or figures (it shows the label), TikZ (export figures to SVG or PNG) and unusual packages.
% Images go in blog/assets/<slug>/ as PNG, SVG or GIF (not PDF). GIFs animate on
% the site but pdflatex can't include them, so a post with GIFs won't compile to PDF.
% Always give images a width with a number in front, e.g. width=0.75\linewidth or
% width=1.0\linewidth. Pandoc ignores a bare width=\linewidth.

\documentclass{article}
\usepackage{amsmath,amssymb,graphicx,booktabs,hyperref}
\graphicspath{{../}} % lets pdflatex find assets/ when compiling from blog/posts/

\newcommand{\E}{\mathbb{E}}

\title{Post Title}
\date{2026-01-01} % must be YYYY-MM-DD

\begin{document}

\begin{abstract}
One or two sentences. This becomes the description on the blog, in the feed and in notification emails.
\end{abstract}

\section{A section}

Text with \emph{emphasis}, inline math $\E[X_i^*]$, and a numbered equation:
\begin{equation}\label{eq:ols}
  \hat\beta = (X^\top X)^{-1} X^\top y.
\end{equation}
Equation~\eqref{eq:ols} can be referred back to.\footnote{Footnotes work too.}

\begin{align*}
  \operatorname{Var}(\hat\beta) &= \sigma^2 (X^\top X)^{-1} \\
                                &\succeq 0.
\end{align*}

\begin{itemize}
  \item A list item.
  \item Another one.
\end{itemize}

% \begin{center}
%   \includegraphics[width=0.75\linewidth,alt={What the figure shows}]{assets/<slug>/figure.png}
% \end{center}

\end{document}
